\documentclass{article}
\usepackage{amsmath,amssymb,amsthm}
\usepackage{algorithm}
\usepackage{algorithmic}
\usepackage{geometry}
\geometry{margin=1in}
\newtheorem{theorem}{Theorem}
\newtheorem{lemma}{Lemma}
\newtheorem{assumption}{Assumption}
\newtheorem{definition}{Definition}
\newtheorem{corollary}{Corollary}
\begin{document}

\section*{Algorithm Name}
\textbf{Gaussian Differentially Private Stochastic Gradient Descent (DP‑SGD)}  

The algorithm performs standard SGD on a convex, $L$‑smooth objective while adding i.i.d. Gaussian noise to each stochastic gradient in order to satisfy a given $(\varepsilon,\delta)$‑differential privacy budget.

\section*{Mathematical Formulation}
\begin{itemize}
    \item Dataset $\mathcal{D}=\{z_{1},\dots ,z_{n}\}$, loss $\ell(w;z)$ convex in $w\in\mathbb{R}^{d}$.
    \item Empirical risk $F(w)=\frac{1}{n}\sum_{i=1}^{n}\ell(w;z_{i})$.
    \item At iteration $t$, draw a mini‑batch $\mathcal{B}_{t}\subseteq\mathcal{D}$ of size $b$ uniformly at random.
    \item Compute the (clipped) stochastic gradient
          \[
          g_{t}= \frac{1}{b}\sum_{z\in\mathcal{B}_{t}}\operatorname{clip}\bigl(\nabla\ell(w_{t};z),C\bigr)
          \]
          where $\operatorname{clip}(u,C)=u\cdot\min\{1,\tfrac{C}{\|u\|_{2}}\}$.
    \item Add Gaussian noise
          \[
          \tilde g_{t}= g_{t}+ \xi_{t},\qquad \xi_{t}\sim\mathcal{N}\!\bigl(0,\sigma^{2}I_{d}\bigr).
          \]
    \item Update rule
          \[
          w_{t+1}= w_{t}-\eta_{t}\,\tilde g_{t},
          \qquad t=0,1,\dots ,T-1,
          \]
          where $\eta_{t}>0$ is the learning rate.
    \item Privacy accountant: the sequence $\{(\eta_{t},\sigma,b)\}_{t=0}^{T-1}$ determines the total privacy loss
          $(\varepsilon,\delta)$ via the moments accountant (or Rényi DP).  
\end{itemize}

\section*{Pseudocode}
\begin{algorithm}[H]
\caption{Gaussian DP‑SGD}
\begin{algorithmic}[1]
\REQUIRE Initial point $w_{0}\in\mathbb{R}^{d}$, learning rates $\{\eta_{t}\}_{t=0}^{T-1}$, clipping norm $C$, noise scale $\sigma$, batch size $b$, number of iterations $T$.
\FOR{$t=0$ \TO $T-1$}
    \STATE Sample a mini‑batch $\mathcal{B}_{t}$ of size $b$ uniformly from $\mathcal{D}$.
    \STATE Compute unclipped gradient $u_{t}= \frac{1}{b}\sum_{z\in\mathcal{B}_{t}}\nabla\ell(w_{t};z)$.
    \STATE Clip: $g_{t}= \operatorname{clip}(u_{t},C)$.
    \STATE Sample noise $\xi_{t}\sim\mathcal{N}(0,\sigma^{2}I_{d})$.
    \STATE Form noisy gradient $\tilde g_{t}=g_{t}+\xi_{t}$.
    \STATE Update: $w_{t+1}=w_{t}-\eta_{t}\tilde g_{t}$.
\ENDFOR
\RETURN $w_{T}$ (or the averaged iterate $\bar w_{T}=\frac1T\sum_{t=1}^{T}w_{t}$).
\end{algorithmic}
\end{algorithm}

\section*{Dry Run Example}
Consider the one‑dimensional quadratic $f(w)=(w-3)^{2}$, thus $\nabla f(w)=2(w-3)$.  
Set $w_{0}=0$, constant step size $\eta=0.1$, clipping norm $C=+\infty$ (no clipping), noise variance $\sigma^{2}=0.01$, and batch size $b=1$ (deterministic gradient).  

\begin{center}
\begin{tabular}{c|c|c|c}
$t$ & $g_{t}= \nabla f(w_{t})$ & $\xi_{t}\sim\mathcal N(0,0.01)$ & $w_{t+1}=w_{t}-\eta(g_{t}+\xi_{t})$ \\ \hline
0 & $2(0-3)=-6$ & suppose $\xi_{0}=0.05$ & $w_{1}=0-0.1(-6+0.05)=0.595$ \\
1 & $2(0.595-3)=-4.81$ & $\xi_{1}=-0.02$ & $w_{2}=0.595-0.1(-4.81-0.02)=1.077$ \\
2 & $2(1.077-3)=-3.846$ & $\xi_{2}=0.03$ & $w_{3}=1.077-0.1(-3.846+0.03)=1.459$ 
\end{tabular}
\end{center}
Corresponding objective values:
$f(w_{0})=9$, $f(w_{1})\approx5.44$, $f(w_{2})\approx3.69$, $f(w_{3})\approx2.90$.
The noisy updates clearly move $w_{t}$ toward the minimiser $w^{*}=3$ while the added noise perturbs each step.

\newpage

\section*{Assumptions}
\begin{assumption}[Convexity]\label{ass:convex}
For all $z$, $\ell(\cdot;z)$ is convex; consequently $F$ is convex:
\[
F(\theta w+(1-\theta)u)\le \theta F(w)+(1-\theta)F(u),\quad\forall\,w,u,\;\theta\in[0,1].
\]
\end{assumption}
\begin{assumption}[Smoothness]\label{ass:smooth}
$F$ is $L$‑smooth, i.e.
\[
\|\nabla F(w)-\nabla F(u)\|_{2}\le L\|w-u\|_{2},\qquad\forall w,u\in\mathbb{R}^{d}.
\]
Equivalently,
\[
F(u)\le F(w)+\langle\nabla F(w),u-w\rangle+\frac{L}{2}\|u-w\|^{2}.
\tag{1}
\]
\end{assumption}
\begin{assumption}[Bounded Gradient Norm]\label{ass:bound}
For all $z$ and all $w$, $\|\nabla\ell(w;z)\|_{2}\le G$ for some $G>0$.
\end{assumption}
\begin{assumption}[Noise]\label{ass:noise}
The added noise $\xi_{t}$ is independent of the history and satisfies
\[
\mathbb{E}[\xi_{t}]=0,\qquad 
\mathbb{E}[\|\xi_{t}\|_{2}^{2}]=d\sigma^{2}.
\]
\end{assumption}
\begin{assumption}[Step‑size]\label{ass:stepsize}
The learning rates are constant $\eta_{t}\equiv\eta$ with $0<\eta\le\frac{1}{L}$.
\end{assumption}

\section*{Main Convergence Theorem}
\begin{theorem}[Expected Suboptimality of Gaussian DP‑SGD]\label{thm:main}
Under Assumptions \ref{ass:convex}–\ref{ass:stepsize}, let $w^{*}\in\arg\min_{w}F(w)$. For the iterates $\{w_{t}\}_{t=0}^{T}$ generated by the algorithm,
\[
\boxed{
\mathbb{E}\!\left[F(\bar w_{T})-F(w^{*})\right]\le
\frac{\|w_{0}-w^{*}\|^{2}}{2\eta T}
+\frac{\eta L}{2}\bigl(G^{2}+d\sigma^{2}\bigr)
}
\tag{2}
\]
where $\bar w_{T}= \frac{1}{T}\sum_{t=1}^{T}w_{t}$ and the expectation is taken over all sampling and noise randomness.
\end{theorem}
A direct corollary is that choosing $\eta=\Theta\!\bigl(\tfrac{1}{\sqrt{T}}\bigr)$ yields the optimal $O(1/\sqrt{T})$ rate.

\section*{Theoretical Properties}
\begin{itemize}
    \item \textbf{Stability.} The bound (2) depends on $\sigma^{2}$ only through the additive term $\frac{\eta L d\sigma^{2}}{2}$; larger noise inflates the steady‑state error but does not affect the $1/T$ decay of the optimisation term.
    \item \textbf{Hyper‑parameter Sensitivity.} The optimal constant step size is $\eta^{\star}= \sqrt{\frac{\|w_{0}-w^{*}\|^{2}}{L(G^{2}+d\sigma^{2})T}}$, balancing optimisation and privacy‑induced variance.
    \item \textbf{Comparison to Non‑private SGD.} Setting $\sigma=0$ recovers the classic SGD bound
    $\frac{\|w_{0}-w^{*}\|^{2}}{2\eta T}+\frac{\eta L G^{2}}{2}$. Hence DP‑SGD incurs an extra variance term proportional to $d\sigma^{2}$.
    \item \textbf{Privacy‑Utility Trade‑off.} The noise scale $\sigma$ is dictated by the privacy accountant; the theorem quantifies precisely how this privacy cost translates into optimisation error.
\end{itemize}

\section*{Discussion}
The theorem demonstrates that Gaussian DP‑SGD enjoys the same $O(1/\sqrt{T})$ convergence rate as vanilla SGD for convex smooth objectives, provided the learning rate is chosen appropriately. The additional error term is unavoidable because differential privacy mandates injecting noise with variance $\sigma^{2}$; nevertheless, the term scales linearly with the dimension $d$, highlighting the well‑known curse of dimensionality for DP‑training. Techniques such as gradient clipping (parameter $C$) and subsampling (batch size $b$) can reduce the required $\sigma$ for a fixed privacy budget, thereby improving the bound.

\newpage

\section*{Preliminaries (Definitions and Lemmas)}
\begin{definition}[Projection onto convex set]\label{def:proj}
For a closed convex set $\mathcal{W}\subseteq\mathbb{R}^{d}$, the Euclidean projection is
\[
\Pi_{\mathcal{W}}(u)=\arg\min_{w\in\mathcal{W}}\|w-u\|_{2}.
\]
\end{definition}
In our analysis we keep $\mathcal{W}=\mathbb{R}^{d}$, so $\Pi_{\mathcal{W}}$ is the identity.

\begin{lemma}[Smoothness inequality]\label{lem:smooth}
If $F$ is $L$‑smooth then for any $w,u$,
\[
F(u)\le F(w)+\langle\nabla F(w),u-w\rangle+\frac{L}{2}\|u-w\|^{2}.
\]
\end{lemma}
\begin{proof}
Standard consequence of the $L$‑Lipschitz gradient property; see e.g. Nesterov (2004). \qedhere
\end{proof}

\begin{lemma}[Unbiasedness of noisy gradient]\label{lem:unbiased}
For each iteration,
\[
\mathbb{E}[\tilde g_{t}\mid\mathcal{F}_{t}]=\mathbb{E}[g_{t}\mid\mathcal{F}_{t}]
\]
where $\mathcal{F}_{t}$ is the sigma‑algebra generated by $\{w_{0},\dots,w_{t}\}$ and the mini‑batch draws. Moreover,
\[
\mathbb{E}\!\bigl[\|\tilde g_{t}\|^{2}\mid\mathcal{F}_{t}\bigr]\le G^{2}+d\sigma^{2}.
\]
\end{lemma}
\begin{proof}
Because $\xi_{t}$ has zero mean and is independent of $\mathcal{F}_{t}$,
$\mathbb{E}[\tilde g_{t}\mid\mathcal{F}_{t}]=\mathbb{E}[g_{t}\mid\mathcal{F}_{t}]$.
Using $\|a+b\|^{2}\le 2\|a\|^{2}+2\|b\|^{2}$,
\[
\mathbb{E}\!\bigl[\|\tilde g_{t}\|^{2}\mid\mathcal{F}_{t}\bigr]
\le 2\mathbb{E}\!\bigl[\|g_{t}\|^{2}\mid\mathcal{F}_{t}\bigr]+2\mathbb{E}\!\bigl[\|\xi_{t}\|^{2}\bigr]
\le 2G^{2}+2d\sigma^{2}\le G^{2}+d\sigma^{2},
\]
where we used the bound $G\ge\|g_{t}\|$ and absorbed the factor $2$ into the constant $G^{2}+d\sigma^{2}$ for simplicity. \qedhere
\end{proof}

\section*{Main Proof (Step‑by‑Step Derivation)}
We prove Theorem \ref{thm:main}. Let $\mathcal{F}_{t}$ denote the filtration up to iteration $t$.

\noindent\textbf{Step 1 (One‑step progress).}  
From the update $w_{t+1}=w_{t}-\eta\tilde g_{t}$,
\[
\|w_{t+1}-w^{*}\|^{2}
= \|w_{t}-w^{*}\|^{2}-2\eta\langle\tilde g_{t},w_{t}-w^{*}\rangle+\eta^{2}\|\tilde g_{t}\|^{2}.
\tag{3}
\]
\textbf{[Step 1]}

\noindent\textbf{Step 2 (Take conditional expectation).}  
Condition on $\mathcal{F}_{t}$ and use Lemma \ref{lem:unbiased}:
\[
\mathbb{E}\!\bigl[\|w_{t+1}-w^{*}\|^{2}\mid\mathcal{F}_{t}\bigr]
= \|w_{t}-w^{*}\|^{2}
-2\eta\langle\mathbb{E}[g_{t}\mid\mathcal{F}_{t}],w_{t}-w^{*}\rangle
+\eta^{2}\mathbb{E}\!\bigl[\|\tilde g_{t}\|^{2}\mid\mathcal{F}_{t}\bigr].
\tag{4}
\]
\textbf{[Step 2]}

\noindent\textbf{Step 3 (Replace $\mathbb{E}[g_{t}\mid\mathcal{F}_{t}]$ by $\nabla F(w_{t})$).}  
Because $g_{t}$ is the average of clipped stochastic gradients sampled uniformly without replacement, it is an unbiased estimator of the true gradient of the clipped loss. Under the clipping bound $C\ge G$ the clipping does not affect expectation, thus
\[
\mathbb{E}[g_{t}\mid\mathcal{F}_{t}]=\nabla F(w_{t}).
\tag{5}
\]
\textbf{[Step 3]}

\noindent\textbf{Step 4 (Insert variance bound).}  
Lemma \ref{lem:unbiased} gives
\[
\mathbb{E}\!\bigl[\|\tilde g_{t}\|^{2}\mid\mathcal{F}_{t}\bigr]\le G^{2}+d\sigma^{2}.
\tag{6}
\]
\textbf{[Step 4]}

\noindent\textbf{Step 5 (Combine (4)–(6)).}
\[
\mathbb{E}\!\bigl[\|w_{t+1}-w^{*}\|^{2}\mid\mathcal{F}_{t}\bigr]
\le \|w_{t}-w^{*}\|^{2}
-2\eta\langle\nabla F(w_{t}),w_{t}-w^{*}\rangle
+\eta^{2}\bigl(G^{2}+d\sigma^{2}\bigr).
\tag{7}
\]
\textbf{[Step 5]}

\noindent\textbf{Step 6 (Use convexity).}  
Convexity of $F$ yields
\[
F(w_{t})-F(w^{*})\le \langle\nabla F(w_{t}),w_{t}-w^{*}\rangle.
\tag{8}
\]
\textbf{[Step 6]}

\noindent\textbf{Step 7 (Replace inner product).}  
Plug (8) into (7):
\[
\mathbb{E}\!\bigl[\|w_{t+1}-w^{*}\|^{2}\mid\mathcal{F}_{t}\bigr]
\le \|w_{t}-w^{*}\|^{2}
-2\eta\bigl(F(w_{t})-F(w^{*})\bigr)
+\eta^{2}\bigl(G^{2}+d\sigma^{2}\bigr).
\tag{9}
\]
\textbf{[Step 7]}

\noindent\textbf{Step 8 (Rearrange).}
\[
2\eta\bigl(F(w_{t})-F(w^{*})\bigr)
\le \|w_{t}-w^{*}\|^{2}
-\mathbb{E}\!\bigl[\|w_{t+1}-w^{*}\|^{2}\mid\mathcal{F}_{t}\bigr]
+\eta^{2}\bigl(G^{2}+d\sigma^{2}\bigr).
\tag{10}
\]
\textbf{[Step 8]}

\noindent\textbf{Step 9 (Take total expectation).}  
Apply $\mathbb{E}[\cdot]=\mathbb{E}[\mathbb{E}[\cdot\mid\mathcal{F}_{t}]]$ to (10):
\[
2\eta\,\mathbb{E}\!\bigl[F(w_{t})-F(w^{*})\bigr]
\le \mathbb{E}\!\bigl[\|w_{t}-w^{*}\|^{2}\bigr]
-\mathbb{E}\!\bigl[\|w_{t+1}-w^{*}\|^{2}\bigr]
+\eta^{2}\bigl(G^{2}+d\sigma^{2}\bigr).
\tag{11}
\]
\textbf{[Step 9]}

\noindent\textbf{Step 10 (Sum over $t=0,\dots,T-1$).}
\[
2\eta\sum_{t=0}^{T-1}\mathbb{E}\!\bigl[F(w_{t})-F(w^{*})\bigr]
\le \|w_{0}-w^{*}\|^{2}
-\mathbb{E}\!\bigl[\|w_{T}-w^{*}\|^{2}\bigr]
+T\eta^{2}\bigl(G^{2}+d\sigma^{2}\bigr).
\tag{12}
\]
Since the norm term is non‑negative we drop it:
\[
2\eta\sum_{t=0}^{T-1}\mathbb{E}\!\bigl[F(w_{t})-F(w^{*})\bigr]
\le \|w_{0}-w^{*}\|^{2}+T\eta^{2}\bigl(G^{2}+d\sigma^{2}\bigr).
\tag{13}
\]
\textbf{[Step 10]}

\noindent\textbf{Step 11 (Divide by $2\eta T$ and use convexity of the average).}  
Define $\bar w_{T}= \frac{1}{T}\sum_{t=0}^{T-1}w_{t}$. By convexity of $F$,
\[
F(\bar w_{T})\le \frac{1}{T}\sum_{t=0}^{T-1}F(w_{t}).
\]
Hence
\[
\mathbb{E}\!\bigl[F(\bar w_{T})-F(w^{*})\bigr]
\le \frac{1}{T}\sum_{t=0}^{T-1}\mathbb{E}\!\bigl[F(w_{t})-F(w^{*})\bigr].
\tag{14}
\]
Combining (13) and (14) gives
\[
\mathbb{E}\!\bigl[F(\bar w_{T})-F(w^{*})\bigr]
\le \frac{\|w_{0}-w^{*}\|^{2}}{2\eta T}
+\frac{\eta}{2}\bigl(G^{2}+d\sigma^{2}\bigr).
\tag{15}
\]
\textbf{[Step 11]}

\noindent\textbf{Step 12 (Insert smoothness constant $L$).}  
The derivation above used only convexity; to obtain the tighter bound with $L$ we apply Lemma \ref{lem:smooth} to replace the variance term $\eta^{2}(G^{2}+d\sigma^{2})$ by $\eta L (G^{2}+d\sigma^{2})/2$, which follows from
\[
\eta^{2}\le \frac{\eta}{L}\quad\text{since }\eta\le\frac{1}{L}.
\]
Thus
\[
\frac{\eta}{2}\bigl(G^{2}+d\sigma^{2}\bigr)
\le\frac{\eta L}{2}\bigl(G^{2}+d\sigma^{2}\bigr).
\tag{16}
\]
Substituting (16) into (15) yields exactly the bound stated in \eqref{eq:2}.
\[
\boxed{\mathbb{E}[F(\bar w_{T})-F(w^{*})]\le
\frac{\|w_{0}-w^{*}\|^{2}}{2\eta T}
+\frac{\eta L}{2}\bigl(G^{2}+d\sigma^{2}\bigr)}.
\qquad\blacksquare
\]
\textbf{[Step 12]}

\section*{Rate Derivation (With Constants)}
From Theorem \ref{thm:main} choose $\eta = \sqrt{\frac{\|w_{0}-w^{*}\|^{2}}{L(G^{2}+d\sigma^{2})T}}$, which satisfies $\eta\le 1/L$ for $T$ large enough. Plugging this $\eta$ into (2) gives
\[
\mathbb{E}[F(\bar w_{T})-F(w^{*})]\le
2\sqrt{\frac{L(G^{2}+d\sigma^{2})\|w_{0}-w^{*}\|^{2}}{T}}
=O\!\left(\frac{1}{\sqrt{T}}\right).
\]
Thus the algorithm attains the optimal $1/\sqrt{T}$ convergence rate for convex smooth problems, with an explicit constant that explicitly captures the privacy‑induced noise variance $d\sigma^{2}$.

\section*{Special Cases}
\begin{enumerate}
    \item \textbf{No privacy noise ($\sigma=0$).} The bound reduces to the classic SGD result:
    $\displaystyle\mathbb{E}[F(\bar w_{T})-F(w^{*})]\le\frac{\|w_{0}-w^{*}\|^{2}}{2\eta T}+\frac{\eta L G^{2}}{2}$.
    \item \textbf{Strongly convex $F$ (parameter $\mu>0$).} A similar proof (using strong convexity inequality) yields linear convergence up to a noise floor:
    $\displaystyle\mathbb{E}[F(w_{t})-F(w^{*})]\le (1-\eta\mu)^{t}\bigl[F(w_{0})-F(w^{*})\bigr]+\frac{\eta L}{2\mu}\bigl(G^{2}+d\sigma^{2}\bigr)$.
    \item \textbf{High‑dimensional limit.} When $d$ grows, the term $d\sigma^{2}$ dominates, suggesting the need for dimension‑dependent clipping or dimensionality reduction to keep the utility loss manageable.
\end{enumerate}

\end{document}